\def\stat{kondrashev}





\def\tit{АРХИТЕКТУРА СИСТЕМЫ ПРЕДОСТАВЛЕНИЯ СЕРВИСОВ ЦИФРОВОЙ ПЛАТФОРМЫ 


ДЛЯ~НАУЧНЫХ ИССЛЕДОВАНИЙ$^*$}





\def\titkol{Архитектура системы предоставления сервисов цифровой платформы 


для %~научных 


исследований}








\def\aut{В.\,А. Кондрашев$^1$}





\def\autkol{В.\,А. Кондрашев}





\titel{\tit}{\aut}{\autkol}{\titkol}





\index{Кондрашев В.\,А.}


\index{Kondrashev V.\,A.}











{\renewcommand{\thefootnote}{\fnsymbol{footnote}}


\footnotetext[1]


{Работа выполнена при частичной поддержке РФФИ (проект 18-29-03091).}}





\renewcommand{\thefootnote}{\arabic{footnote}}


\footnotetext[1]{Институт проблем информатики Федерального исследовательского центра 


<<Информатика и~управление>> Российской академии наук, 


\mbox{VKondrashev@frccsc.ru}}








 \label{st\stat}





                  \thispagestyle{headings}


                  


  \vspace*{-12pt}





  





      \Abst{Рассматриваются вопросы предоставления научных услуг в~условиях 


цифровой экономики на основе цифровой платформы для научных исследований. 


Предложена биз\-нес-модель научных сервисов цифровой платформы. Делается акцент на 


возможностях цифровой платформы как инструмента для управления научными 


исследованиями в~цифровой экономике. Описываются основные архитектурные решения 


системы предоставления сервисов цифровой платформы для научных исследований.}


      


      \KW{цифровая платформа; облачный сервис; научный сервис; информационные 


технологии; информационная поддержка де\-я\-тель\-ности; управление научными сервисами; 


единая информационная среда}





\DOI{10.14357\!/\!08696527180310} 








\vskip 10pt plus 9pt minus 6pt








      \thispagestyle{headings}


      


      \vspace*{-18pt} 


      


\section{Введение}





\vspace*{-3pt}





     В настоящее время в~соответствии с~программой <<Цифровая 


экономика Российской Федерации>>~[1] и~Указом Президента РФ 


<<О~национальных целях и~стратегических задачах развития Российской 


Федерации на период до~2024~года>>~[2] начаты работы по созданию 


отраслевых цифровых платформ. Одно из базовых направлений  


программы~--- создание исследовательской инфраструктуры. Очевидно, для 


повышения эффективности ее использования в~интересах научных 


коллективов России требуется создание цифровой платформы для научных 


исследований. 


     


     Под цифровыми платформами понимают общее информационное 


пространство поставщиков и~потребителей товара (услуги), 


предоставляющее биз\-нес-мо\-дели взаимодействия на основе  


информационных тех\-но\-ло\-гий. Заметим, что\linebreak циф\-ровая платформа позволяет 


накапливать информацию о взаимодействии поставщиков и~потребителей, 


качестве оказанных услуг, снижать число транзакционных связей и~наделять 


биз\-нес-мо\-де\-ли взаимодействия поставщика и~потребителя элементами 


искусственного интеллекта. В~рамках работ по программе развития 


цифровой экономики в~России весной~2018~г.\ принято следующее 


определение цифровой платформы~[3]: <<Цифровая платформа~--- это 


система алгоритмизированных взаимовыгодных взаимоотношений 


значимого количества независимых участников отрасли экономики (или 


сферы деятельности), осуществляемых в~единой информационной среде, 


приводящая к~снижению транзакционных издержек за счет применения 


пакета цифровых технологий работы с~данными и~изменения системы 


разделения труда>>.


     


     Бизнес-процессы проведения научных исследований обладают рядом 


особенностей, которые должны быть учтены при создании цифровой 


платформы. Так, в~[4] цифровая платформа представлена в~составе трех 


компонентов: центра компетенции, центра обработки данных и~комплекса 


научных сервисов. При этом понятие <<научный сервис>> является 


ключевым для цифровой платформы для научных исследований 


и~определяется как совокупность действий (процессов) и~средств 


обеспечения процесса (ресурсов) по обслуживанию выполнения работ 


научно-исследовательского и~прикладного характера путем предоставления 


потребителю оборудования, расходных материалов, ин\-фор\-ма\-ци\-он\-но-ком\-му\-ни\-ка\-ци\-он\-ных 


ресурсов, обеспечивающих ресурсов, продуктов 


интеллектуальной научной деятельности и~обслуживающих человеческих 


ресурсов (субъекты сервисной деятельности), результатом которых является 


научная услуга. Отметим также, что в~статье понятие <<бизнес>> трактуется 


более широко, чем его классическое определение: <<экономическая, 


предпринимательская деятельность, направленная на извлечение прибыли>>. 


Отнесем к~этому понятию в~рамках настоящей статьи также  


на\-уч\-но-тех\-ни\-че\-скую деятельность, направленную на получение 


нового знания (результата интеллектуальной деятельности), в~том числе при 


употреблении в~сочетаниях бизнес-процессы, биз\-нес-мо\-дели.


     


\section{Бизнес-модели предоставления научных сервисов цифровой~платформы}





     Как показал анализ биз\-нес-мо\-де\-лей в~области инноваций, наиболее 


близок к~концепции цифровой платформы~[5] подход~[6], представляющий  


биз\-нес-мо\-дель инновации как комплекс следующих элементов, 


обеспечивающих в~совокупности создание и~распространение новой 


ценности:


     \begin{itemize}


\item потребляемая ценность;\\[-13.5pt]


     \item  ключевые ресурсы (персонал, оборудование и~др.);\\[-13.5pt]


     \item  ключевые процессы (показатели и~критерии, регламенты, 


протоколы и~др.);\\[-13.5pt]


     \item формула прибыли.


     \end{itemize}


     


     Поясним эти элементы модели с~точки зрения научного исследования, 


в~котором потребляемой ценностью является научный сервис (научная 


услуга). 


     


     Описание потребляемой ценности содержит понимание того, что 


требуется от научного сервиса с~точки зрения его потребителя. Оно 


включает:


     \begin{itemize}


     \item  характеристики ключевого потребителя исследования;\\[-13.5pt]


     \item спецификацию потребности, требующей удовлетворения;\\[-13.5pt]


\item  предложение исследовательского коллектива по оказанию научной 


услуги.


\end{itemize}





     Ключевые ресурсы характеризуют персонал, используемые 


технологии, обору\-до\-ва\-ние, информационные ресурсы, каналы 


взаимодействия с~партнерами и~потребителями и~пр.


     


     Ключевые процессы дают понимание того, как организованы 


исследования, чтобы можно было постоянно и~в~нарастающем объеме 


предлагать научные услуги. Обычно в~этот элемент включаются также 


система функциональных показателей, правила и~нормы организации 


процессов.


     


     Формула прибыли описывает модель получения результата, структуру 


затрат, скорость обращения (возобновления) ресурсов.


     


     Соответственно, схема действий при формировании научных сервисов 


циф\-ро\-вой платформы будет выглядеть следующим образом (рис.~1): 


     \begin{enumerate}[(1)]


\item определение потребительской ценности научного сервиса, основываясь 


на реальных потребностях пользователей услуги;


     \item составление формулы прибыли: определение источников 


финансирования, решения по поводу того, какие ресурсы и~процессы 


потребуются для выполнения научного сервиса;


     \item сравнение новой модели с~уже существующей, решение о том, 


можно ли оказать научную услугу в~существующей организационной 


структуре или потребуется создавать новое подразделение.


     \end{enumerate}


          \begin{figure} %fig1


     \vspace*{1pt}


 \begin{center}


 \mbox{%


 \epsfxsize=104.106mm 


 \epsfbox{kon-1.eps}


 }


\end{center}


\vspace*{-9pt}


     \Caption{Бизнес-модель научного сервиса цифровой платформы}


      \end{figure}


          При этом надо стремиться, чтобы цифровая платформа была способна 


поддержать биз\-нес-мо\-де\-лей разного объема оказания научного сервиса 


(по классификации~[7]):


     \begin{itemize}


     \item интегрированную модель с~охватом всей цепочки выполнения 


исследования, существенным вкладом в~получение качественного результата и~доступом к~необходимым комплементарным активам внутри организации;


     \item модель <<дирижера>> (Orchestrator), при которой оказание 


услуги специализируется на одном или нескольких участках исследования, 


обладает существенным вкладом в~получение качественного общего 


результата и~доступом к~ключевым комплементарным активам, в~том числе 


за счет сотрудничества с~другими научными коллективами; %\\ %[-14pt]


     \item модель <<игрока, действующего на одном уровне>> (Layer 


Player), при которой оказание услуги специализируется на одном звене 


исследования, имеется относительно небольшой вклад в~общий результат 


исследования, присутствует доступ к~ключевым комплементарным активам 


в~результате сотрудничества с~другими научными коллективами; %\\ %[-14pt]


     \item  модель <<маркетмейкера>> (Market Maker), при которой 


выполняется совершенно новое звено исследования, характеризуется 


относительно небольшим вкладом в~общий результат исследования 


и~высоким спросом на новаторский научный сервис.


     \end{itemize}


     


\section{Аналитика результатов предоставления научных сервисов


цифровой~платформы}





     Цифровая платформа для научных исследований помимо создания 


инновационной основы исследовательским организациям для 


предоставления научных сервисов может обеспечить структуризацию 


научных сервисов, создание сквозных технологий для эффективного 


использования и~учета, а~также систематизацию и~оптимизацию затрат на 


проведение научных исследований.


     


     Целесообразно использовать цифровую платформу как один из 


инструментов управления научными исследованиями. В~цифровой 


платформе можно описать управление научными исследованиями в~виде 


единого биз\-нес-про\-цес\-са~[8], целью которого являются научные 


результаты (результаты интеллектуальной деятельности), а управляющими 


воздействиями~--- ресурсы, направляемые на проведение исследования.


     


     Информация, доступная и~накапливаемая цифровой платформой, 


позволит выполнить аналитическую оценку результатов предоставления 


научных сервисов (проведения научных исследований) с~учетом следующих 


положений, свойственных процессам управления научными исследованиями:


     \begin{itemize}


     \item результаты научного исследования могут переоцениваться во 


времени; %[-14pt]


     \item экспертная оценка результатов исследования носит субъективный 


характер и~должна взвешиваться с~использованием рейтинга эксперта; %[-14pt]


     \item рейтинг эксперта наряду с~наукометрическими показателями 


(ученая степень, звание, публикации и~т.\,д.)\ должен включать показатели 


качества проведенных экспертиз и~выполненных работ.


     \end{itemize}


     


     Накопление в~информационном пространстве цифровой платформы 


информации по экспертам, научным коллективам, научным сервисам 


и~научным результатам обеспечит применение известных методов 


наукометрии для научных исследований, организуемых с~использованием 


цифровой платформы.


     


     Цифровая платформа позволит взвешенно оценить актуальность, 


важность, результативность научных сервисов и~распределить ресурсы 


с~использованием принципов состязательности, целесообразности 


и~оптимальности.


     


\section{Основные архитектурные решения системы предоставления 


научных сервисов цифровой платформы}





     С учетом вышеизложенных аспектов определим состав цифровой 


платформы научных исследований как совокупность следующих 


компонентов:


     \begin{itemize}


     \item система описания научного сервиса;\\[-13.5pt]


     \item система публикации научного сервиса;\\[-13.5pt]


     \item система классификации научных сервисов;\\[-13.5pt]


     \item система поиска научного сервиса;\\[-13.5pt]


     \item система заказа научного сервиса;\\[-13.5pt]


     \item система планирования и~учета ресурсов;\\[-13.5pt]


     \item система учета результатов и~экспертных оценок;\\[-13.5pt]


     \item система доступа пользователей.


     \end{itemize}


     


     Система предоставления научных сервисов цифровой платформы 


создает единое информационное пространство для функционирования 


научных сервисов, смежных систем и~вспомогательных сервисов. Ключевой 


смежной системой для системы предоставления научных сервисов является 


аналитическая система. Совместное функционирование этих двух систем 


в~едином информационном пространстве цифровой платформы обеспечит 


выполнение единого бизнес-процесса управления научными исследованиями 


в~рамках цифровой платформы.


     


     Рассмотрим кратко назначение и~основные функции компонентов 


системы предоставления научных сервисов.


     


     Система описания научного сервиса является основой системы 


предоставления научных сервисов и~предназначена для спецификации 


карточки научного сервиса, определяющей следующие его основные 


характеристики:


     \begin{itemize}


     \item  идентификационные данные;\\[-13.5pt]


     \item данные поставщика сервиса;\\[-13.5pt]


     \item форму заказа сервиса;\\[-13.5pt]


     \item технологическую карту (методику, процесс) выполнения 


исследования с~указанием потребляемых ресурсов (трудозатраты персонала, 


время работы приборов, потребление расходных материалов и~т.\,п.);\\[-13.5pt]


     \item процедуры заказа сервиса и~согласования изменений 


в~технологической карте (методике) исследования.


     \end{itemize}


     


     Система публикации научного сервиса ведет каталог доступных для 


заказа научных сервисов, обеспечивает систематизацию <<близких>> 


научных сервисов.


     


     Система классификации научных сервисов предназначена для 


группировки научных сервисов в~соответствии с~заданными критериями 


(область исследования, методики исследования, приборная база, 


территориальная и~организационная принадлежность и~т.\,д.)\ с~целью 


обеспечения быстрого релевантного поиска и~систематизации научных 


сервисов. Для системы классификации предусматривается изменяемый 


список критериев с~динамически перестраиваемой системой группировки.


     


     Система поиска научного сервиса предназначена для нахождения 


<<близких>> научных сервисов с~использованием системы классификации 


научных сервисов и~данных систем описания научного сервиса, учета 


результатов и~экспертных оценок. В~системе предусматриваются блоки 


поиска <<близких>> сервисов с~использованием методов искусственного 


интеллекта с~элементами самообучения.


     


     Система заказа научного сервиса предназначена для оформления 


договорных отношений на оказание научной услуги в~ходе выполнения 


процесса выбора и~заказа научной услуги во взаимодействии с~системами 


публикации, поиска, описания научного сервиса, планирования и~учета 


ресурсов.


     


     Система планирования и~учета ресурсов предназначена для ведения 


реестра ресурсов (персонала, приборной базы, расходных материалов) и~их 


состояния, а~так\-же календаря загрузки и~резервирования ресурсов.


     


     Система учета результатов и~экспертных оценок предназначена для 


ведения реестра договоров, результатов интеллектуальной деятельности, 


включая публикации, согласованных изменений методик и~технологических 


карт исследований, результатов экспертной деятельности.


     


     Система доступа пользователей предназначена для ведения реестра 


пользователей, обеспечения интерфейса доступа к~функциям системы 


в~соответствии с~полномочиями пользователей, определяемыми их ролями 


(поставщик научной услуги, оператор оборудования, потребитель научной 


услуги, эксперт, аналитик и~т.\,д.). В~системе предусматривается 


изменяемый перечень ролей.


     


     Для создания единого информационного пространства цифровой 


платформы с~высокой степенью масштабирования, эластичности 


и~адаптируемости целесообразно ориентироваться на технологии 


виртуализации (контейнеризации) вычислительных ресурсов, облачных 


вычислений, микросервисных архитектур и~интеграционных шин.


     


     В этом случае с~учетом подходов~[9, 10] основные архитектурные 


решения системы предоставления научных сервисов формируются как 


сер\-вис\-но-ори\-ен\-ти\-ро\-ван\-ный облачный комплекс систем микросервисной 


архитектуры, связан\-ных интеграционной шиной, предоставляющий услуги 


пользователям с~использованием веб-портальных технологий (рис.~2).


     








\section{Заключение}





\begin{figure} %fig2


 \vspace*{1pt}


 \begin{center}


 \mbox{%


 \epsfxsize=128.996mm 


 \epsfbox{kon-2.eps}


 }


\end{center}


\vspace*{-6pt}


\Caption{Система предоставления научных сервисов цифровой платформы научных 


исследований}


\end{figure}








     Представленные в~статье бизнес-модели организации научных 


исследований как сервисов цифровой платформы, основные архитектурные 


решения системы предоставления сервисов цифровой платформы для 


научных исследований служат современным подходом к~организации 


и~управлению научными исследованиями в~условиях цифровой экономики.


     


     Представленные подходы к~созданию цифровой платформы научных 


исследований развиваются в~ФИЦ ИУ РАН в~рамках государственного 


задания и~грантов.





%\vspace*{-9pt}


     


{\small\frenchspacing


{%\baselineskip=10.8pt


\addcontentsline{toc}{section}{Литература}


\begin{thebibliography}{99}


\bibitem{1-ksh}


Цифровая экономика Российской Федерации: Программа, утверж\-ден\-ная Распоряжением 


Правительства РФ от 28.07.2017 №\,1632-р. {\sf 


http://static.government.\linebreak ru/media/files/9gFM4FHj4PsB79I5v7yLVuPgu4bvR7M0.pdf}.


\bibitem{2-ksh}


О~национальных целях и~стратегических задачах развития Российской Федерации на 


период до 2024 года: Указ Президента РФ от~7~мая 2018~г.\ №\,204. {\sf 


http:// static.kremlin.ru/media/acts/files/0001201805070038.pdf.}


\bibitem{3-ksh}


Цифровые платформы: подходы к~определению и~типизации. Ростелеком. {\sf  


http://\linebreak d-russia.ru/wp-content/uploads/2018/04/digital\_platforms.pdf.}


\bibitem{4-ksh}


\Au{Зацаринный А.\,А., Горшенин А.\,К., Волович К.\,И., Колин~К.\,К., Кондрашев~В.\,А., 


Степанов~П.\,В.} Управление научными сервисами как основа национальной цифровой 


платформы <<Наука и~образование>>~/\!\!/ Стратегические приоритеты, 2017. №\,2(14). 


С.~103--113.


\bibitem{5-ksh}


\Au{Зацаринный А.\,А., Горшенин А.\,К., Волович~К.\,И., Кондрашев~В.\,А.} Основные 


направления развития информационных технологий в~условиях вызовов цифровой 


экономики~/\!\!/ Цифровая обработка сигналов, 2018. №\,1. С.~3--7.


\bibitem{6-ksh}


\Au{Джонсон М., Кристенсен К., Кагерманн~Х.} Обновление биз\-нес-мо\-де\-ли~/\!\!/ 


Harvard Bus. Rev., 2009. {\sf https://hbr-russia.ru/management/strategiya/a9780}.


\bibitem{7-ksh}


\Au{Швайцер Л.} Концепция и~эволюция бизнес-моделей~/\!\!/ ЭКОВЕСТ, 2007. Т.~6. 


Вып.~2. С.~146--168. {\sf http://www.research.by/webroot/delivery/files/2007n2r01.pdf}.


\bibitem{8-ksh}


Исследование вопросов управления результатами научно-исследовательской деятельности 


организаций, подведомственных ФАНО России, и~научными сервисами сети ЦКП ФАНО: 


Отчет о~НИР <<Сервис-У>>.~--- ФИЦ ИУ РАН, 2016. 437~с.





\bibitem{10-ksh}


\Au{Волович К.\,И., Денисов С.\,А., Кондрашев~В.\,А., Сучков~А.\,П.} Методология 


создания веб-сервисного информационного взаимодействия в~системе распределенных 


ситуационных центров~/\!\!/ Системы и~средства информатики, 2016. Т.~26. №\,4.  


С.~51--59.





\bibitem{9-ksh}


\Au{Волович К.\,И., Зацаринный~А.\,А., Кондрашев~В.\,А., Шабанов~А.\,П.} О~некоторых 


подходах к~представлению научных исследований как облачного сервиса~/\!\!/ Системы 


и~средства информатики, 2017. Т.~27. №\,1. С.~73--84.





        \end{thebibliography}


} }








\vspace*{-3pt}





\hfill{\small\textit{Поступила в~редакцию 18.07.18}}














\vspace*{8pt}





\hrule





\vspace*{2pt}





%\newpage





%\vspace*{-28pt}





\hrule





%\vspace*{9pt}











\def\tit{ARCHITECTURE OF THE SERVICE DELIVERY SYSTEM 


FOR~THE~RESEARCH SERVICES DIGITAL PLATFORM}





\def\titkol{Architecture of~the~service delivery system 


for~the~research services digital 


platform}








\def\aut{V.\,A.~Kondrashev}


\def\autkol{V.\,A.~Kondrashev}











\titel{\tit}{\aut}{\autkol}{\titkol}





\vspace*{-9pt}





\noindent


      Institute of Informatics Problems, Federal Research Center ``Computer 


Science and~Control'' of the Russian Academy of Sciences, 44-2~Vavilov Str., 


Moscow 119333, Russian Federation








 \noindent





\def\leftfootline{\small{\textbf{\thepage}


\hfill Sistemy i Sredstva Informatiki~---


Systems and Means of Informatics 2018 vol~28 no\,3}


}%


 \def\rightfootline{\small{Sistemy i Sredstva Informatiki~---


 Systems and Means of Informatics   2018   vol~28   no\,3


\hfill \textbf{\thepage}}}      


      


           


      \Abste{The paper deals with providing scientific services in 


      the digital economy using a~digital platform for scientific 


      research. The authors propose a~business model of a~research 


      services digital platform. The capabilities of a~digital platform as a~tool 


      for managing scientific research in the digital economy are discussed. 


      The main architectural solutions for the system providing 


      scientific services of the digital platform are described.}


      


      \KWE{digital platform; cloud service; scientific service; 


      information technology; information support activities; 


      research management; unified information environment}


      


\DOI{10.14357\!/\!08696527180310} 





%\vspace*{-24pt}





\Ack


\noindent


The work was partly supported by the Russian Foundation for


Basic Research (project 18-29-03091).





 








\renewcommand{\bibname}{\protect\rmfamily References}


%\renewcommand{\bibname}{\large\protect\rm References}





{\small\frenchspacing


{%\baselineskip=10.8pt


\addcontentsline{toc}{section}{References}


\begin{thebibliography}{99}





\bibitem{1-ksh-1}


Tsifrovaya ekonomika Rossiskoy Federatsii: Programma, utverzhdennaya Rasporyazheniem Pravitel'stva


RF ot 28.07.2017 No.\,1632-r


[The program ``Digital Economy of the Russian Federation'' (approved by the Order of the 


Government of the Russian Federation dated 28.07.2017 No.\,1632-r)]. 2017. Available at: {\sf 


http://\linebreak static.government.ru/media/files/9gFM4FHj4PsB79I5v7yLVuPgu4bvR7M0.pdf} 


(accessed July~20, 2018).


\bibitem{2-ksh-1}


O~natsional'nykh tselyakh i~strategicheskikh zadachakh razvitiya Rossiyskoy


Federatsii na period do 2024~goda: Ukaz Presidenta RF ot 7~maya 2018~g. No.\,204


[Decree of the President of the Russian Federation


dated May~7, 2018, No.\,204 ``On national 


goals and strategic tasks for the development of the Russian Federation for the period up to 


2024'']. 2018. Available at: {\sf 


http://static.kremlin.ru/media/ acts/files/0001201805070038.pdf} (accessed July~20, 2018).


\bibitem{3-ksh-1}


Rostelecom. 2018. 


Tsifrovye platformy: podkhody k~opredeleniyu i~ti\-pi\-za\-tsii


[Digital platforms approaches to definition and typification].


Available at: {\sf  


http://\linebreak d-russia.ru/wp-content/uploads/2018/04/digital\_platforms.pdf} (accessed July~20, 


2018).


\bibitem{4-ksh-1}


\Aue{Zatsarinnyy, A.\,A., A.\,K.~Gorshenin, K.\,I.~Volovich, K.\,K.~Kolin, 


V.\,A.~Kondrashev, and P.\,V.~Stepanov.} 2017. 


Upravlenie nauchnymi servisami kak osnova natsional'noy tsifrovoy platformy


``Nauka i~obrazovanie''


[Management of scientific services as the 


basis of the national digital platform ``Science and Education''].


 \textit{Strategicheskie prioritety}


 [Strategic Priorities] 


2(14):103--113.


\bibitem{5-ksh-1}


\Aue{Zatsarinnyy, A.\,A., A.\,K. Gorshenin, K.\,I.~Volovich, and V.\,A.~Kondrashev}. 


2018.  Osnovnye napravleniya razvitiya informatsionnykh tekhnologiy


v~usloviyakh vyzovov tsifrovoy ekonomiki


[The main directions of the development of information technologies in the face of the 


challenges of the digital economy]. \textit{Digital Signal Processing} 1:3--7.


\bibitem{6-ksh-1}


\Aue{Johnson, M., K.~Christensen, and H.~Kagermann}. 2009. Updating the business 


model. \textit{Harvard Bus. Rev.} Available at: {\sf  


https://hbr-russia.ru/management/\linebreak strategiya/a9780} (accessed July~20, 2018).


\bibitem{7-ksh-1}


\Aue{Schweizer, L.} 2007. The concept and evolution of business models. 


\textit{ECOVEST} 6(2):146--168. Available at: {\sf 


http://www.research.by/webroot/delivery/\linebreak  files/2007n2r01.pdf} (accessed July~20, 2018).


\bibitem{8-ksh-1}


FRC CSC RAS. 2016. 


Issledovanie voprosov upravleniya rezul'tatami nauchno-issledovatel'skoy


deyatel'nosti organizatsiy,


podvedomstvennykh FANO Rossii, i~nauchnymi servisami seti


TsKP FANO:


Otchet o~NIR ``Servis-U''


[Report on the research ``Service-U.'' Study of the issues of 


managing the results of research activities of organizations subordinate to FASO 


of Russia and 


the scientific services of the FASO CKP]. 437~p. 





\bibitem{10-ksh-1}


\Aue{Volovich, K.\,I., S.\,A.~Denisov, V.\,A.~Kondrashev, and A.\,P.~Suchkov}. 2016. 


Me\-to\-do\-lo\-giya sozdaniya veb-servisnogo informatsionnogo vzaimodeystviya v~sisteme 


ras\-pre\-de\-len\-nykh situatsionnykh tsentrov [Methodology of creating web-service interactions 


in the systems of distributed situational centers]. \textit{Sistemy i~Sredstva Informatiki~--- 


System and Means of Informatics} 26(4):51--59.


\bibitem{9-ksh-1}


\Aue{Volovich, K.\,I., A.\,A. Zatsarinnyy, V.\,A.~Kondrashev, and A.\,P.~Shabanov.} 2017.  


O~nekotorykh podkhodakh k~predstavleniyu nauchnykh issledovaniy kak oblachnogo servisa 


[Scientific research as a~cloud service]. \textit{Sistemy i~Sredstva Informatiki~--- System 


and Means of Informatics} 27(1):73--84.


\end{thebibliography}


} }








\vspace*{-3pt}





\hfill{\small\textit{Received July 18, 2018}} 


      


      \Contrl


      


      \noindent


\textbf{Kondrashev Vadim A.} (b.\ 1963)~--- senior scientist, Institute of Informatics 


Problems, Federal Research Center ``Computer Science and Control'' of the Russian Academy of 


Sciences, 44-2~Vavilov Str., Moscow 119333, Russian Federation; 


\mbox{VKondrashev@frccsc.ru}


      


\label{end\stat}








\renewcommand{\bibname}{\protect\rm Литература}


     


